% Integrated F.49 after root, topos and finite-field review.
\subsection{Hilbertian normal extensions and small geometric Galois groups}

\begin{claim}{F.49}{Proved}\label{clm:F49}
Let \(D/\mathbf Q\) be finitely generated of positive transcendence
degree. Fix an algebraic closure containing \(\overline{\mathbf Q}\),
and let \(k/D\) be algebraic Galois with
\(\overline{\mathbf Q}\subset k\).
If \(k\) is Hilbertian, every finitely generated extension \(L/k\)
has closed-point-preserving all-scheme étale-topos reconstruction.

A profinite group is \emph{small} if it has finitely many open
subgroups of each prescribed finite index. If
\(\operatorname{Gal}(k/D\overline{\mathbf Q})\) is small, then \(k\)
is Hilbertian, and every subfield \(k_0\subset L\) of positive
transcendence degree over \(\mathbf Q\), for every finitely generated
\(L/k\), reconstructs. The positive-transcendence restriction is
essential to this statement: \(\overline{\mathbf Q}\) itself fails
by F.41.

This class contains fields outside the entire sufficient class of
F.46, even after changing the finitely generated algebraic model.
Two examples, with \(D=\mathbf Q(t)\), are the following:
\begin{enumerate}
\item Let \(\Omega_m/D\) be the splitting field of \(P_m(X)-t\), where
\[
 P_m(X)=X^{2m+1}-\frac{2(2m+1)}{m+1}X^{m+1}+(2m+1)X,\qquad m\ge2.
 \tag{F.49.1}
\]
Then the compositum
\[
 \begin{gathered}
 k_{\mathrm{alt}}
   =\overline{\mathbf Q}(t)\prod_{m\ge2}\Omega_m,\\
 \operatorname{Gal}(k_{\mathrm{alt}}/\overline{\mathbf Q}(t))
       \simeq\prod_{m\ge2}A_{2m+1}.
 \end{gathered}
 \tag{F.49.2}
\]
\item Let \(k_{\mathrm{tri}}\) be the maximal extension of
\(\overline{\mathbf Q}(t)\) unramified over
\(\mathbf P^1_{\overline{\mathbf Q}}\setminus\{0,1,\infty\}\).
Then
\[
 \operatorname{Gal}(k_{\mathrm{tri}}/\overline{\mathbf Q}(t))
       \simeq\widehat F_2.
 \tag{F.49.3}
\]
This field is not radically closed: it contains all roots of \(t\)
and \(1-t\), but contains no square root of \(t-2\).
\end{enumerate}
Both fields, all their finitely generated extensions, and the
positive-transcendence subfields specified above reconstruct.
No necessity of Hilbertianity or smallness is asserted. In particular,
this does not assert reconstruction for every normal algebraic
extension of a positive-transcendence finitely generated field, and
F.9 remains unresolved.
\end{claim}
\begin{proof}
\proofparagraph{Normal finite-type descent and the divisible kernel}
For a finitely generated \(L/k\) and any finite list of geometric
data, choose a finitely generated field \(F/\mathbf Q\), containing
\(D\), over which the data are defined and such that
\[
 L=Fk,\qquad L/F\ \hbox{is algebraic Galois},\qquad
 B=F\overline{\mathbf Q}\subset L.
 \tag{F.49.4}
\]
Normality follows by base change from \(k/D\).
The Hilbertian core will specialize over \(L\) itself; it does not
assume Hilbert irreducibility over this arithmetic descent field \(F\).

For every finite extension \(E/k\), we claim
\[
 D(E^\times):=\bigcap_{n\ge1}(E^\times)^n\subset k^\times.
 \tag{F.49.5}
\]
Take \(a\in D(E^\times)\). Descending the finitely many equations
defining \(E/k\), and taking a normal closure, gives a finite Galois
extension \(M/D\) such that \(a\in M\) and \(E\subset Mk\).
For \(\sigma\in\operatorname{Gal}(Mk/k)\), the normality of \(M/D\)
makes \(M\) stable under \(\sigma\). For each \(n\), choose
\(z_n\in E\) with \(z_n^n=a\). If
\(\tau\in\operatorname{Gal}(Mk/M)\), then \(\tau\) preserves \(k\),
because \(k/D\) is normal. The automorphisms \(\sigma,\tau\) commute
on \(M\) and on \(k\), hence on \(Mk\). Moreover \(\sigma\) fixes
all roots of unity. Since \(\tau(z_n)/z_n\in\mu_n\), it follows that
\(\sigma(z_n)/z_n\) is fixed by every such \(\tau\), and belongs to
\(M\). Thus
\[
                 \sigma(a)/a\in D(M^\times)=\{1\}.
\]
The last equality is the divisor and number-field unit argument
in F.4 and F.33, applied to the finitely generated field \(M/\mathbf Q\).
Consequently \(a\) is fixed by \(\operatorname{Gal}(Mk/k)\), proving
(F.49.5). No equality \(D(E^\times)=k^\times\) is needed or claimed.

The same inclusion holds for every finitely generated \(E/k\).
Let \(k_1\) be the relative algebraic closure of \(k\) in \(E\);
it is finite over \(k\). On a normal proper model of \(E/k_1\),
an infinitely divisible element has zero valuation at every prime
divisor. It therefore lies in \(k_1^\times\). Each of its \(n\)-th
roots in \(E\) is algebraic over \(k_1\), and hence is also in \(k_1\),
by relative algebraic closedness. The element belongs to
\(D(k_1^\times)\), so the finite-extension argument applies.
This proves the precise scalar-kernel inclusion that will be used
in point evaluation.

\proofparagraph{Hilbertian finite-quotient detection}
If \(k\) is Hilbertian, every finitely generated \(L/k\) is Hilbertian.
Finite extensions preserve Hilbertianity; a finitely generated
transcendental extension of any field is Hilbertian by
\citet[Theorem~1.2 and Corollary~3.2]{F_BarySorokerParan2013}.
The finite-extension and full-degree specialization statements are
also those used in F.44--F.45
\citep[Proposition~3.3.1 and Theorem~3.4.1]{F_Serre2007}.
In particular \(L\) and every finite extension of \(L\) have
nontrivial finite separable extensions.

Let \(X/L\) be smooth and geometrically connected, and let
\(Y\to X\) be a connected finite étale Galois cover with group \(G\).
It represents an arbitrary finite quotient of the arithmetic
\(\pi_1(X)\). Suppose first that \(X\) is a curve.
On its smooth compactification, Riemann--Roch gives rational
functions \(X\to\mathbf P^1\) of arbitrarily large degree \(d\).
After shrinking, the function and the composite
\(Y\to\mathbf P^1\) are finite étale. The latter connected cover
has degree \(d|G|\), including its arithmetic constant extension.
Hilbert irreducibility gives an \(L\)-rational specialization whose
entire fibre is one field of that degree. Its \(G\)-quotient is a
closed point \(x\in X\) of degree \(d\), and the fibre over \(x\)
is a field of degree \(|G|\) over \(L(x)\). Therefore
\[
                      G_{L(x)}\longrightarrow G
                      \quad\hbox{is surjective}.
 \tag{F.49.6}
\]
Removing a prescribed finite set before specialization and
increasing \(d\) gives detecting points of unbounded degree in
every nonempty open. For higher-dimensional \(X\), F.29's
cover-preserving smooth-curve construction reduces to this curve
case, retaining the arithmetic connectedness and any prescribed
open restriction. This proves exactly the finite-quotient
detection used in F.45--F.46, including the constant quotient.

\proofparagraph{Integral Kummer residues in one fixed norm field}
Let
\[
 U=\operatorname{Spec}L[T,1/(P_1\cdots P_r)]
\]
for distinct monic irreducible polynomials \(P_j\), and let
\(x_0\in U(L)\). The normalized compatible Kummer class has the form
\[
 \kappa(b_0)+
   \sum_j\lambda_j\bigl([P_j(T)]-[P_j(x_0)]\bigr),
       \qquad\lambda_j\in\widehat{\mathbf Z}.
 \tag{F.49.7}
\]
Assume that its evaluation at every point over every finite extension
of \(L\) is represented by an actual nonzero field element.
This is the property supplied by the pinned morphisms in question.

Choose \(F\) in (F.49.4) containing the coefficients and normalizing
data. Choose a nontrivial finite extension of \(L\), take its
Galois closure, and descend its defining equations to a sufficiently
large such \(F\). Take a finite Galois extension
\(M/\mathbf Q(t_1,\ldots,t_s)\) containing \(F\), all the roots of
the \(P_j\), and these descended extension data.
Then \(Mk/L\) is finite Galois and nontrivial.
Choose a nonidentity
\(\sigma\in\operatorname{Gal}(Mk/L)\).
It preserves \(M\), fixes \(k\) and \(\overline{\mathbf Q}\),
and acts nontrivially on \(M\).

Apply F.45's moved-geometric-divisor construction to \(M\) over
its algebraic constants \(M\cap\overline{\mathbf Q}\).
It gives pairwise distinct geometric discrete valuations \(v_i\),
with disjoint \(\sigma\)-translates, at which all prescribed
coefficients and distinct-root differences are units.
Semilocal approximation, in the localization at these
height-one primes, gives \(\alpha\in U(M)\) such that
\[
 v_i(P_j(\alpha))=\delta_{ij},\qquad
 v_i(P_j(\sigma\alpha))=0.
 \tag{F.49.8}
\]
The valuation-one assertion is the root-plus-uniformizer Taylor
calculation in F.45; \(\sigma\) fixes every \(P_j\), but may permute
its roots. Put \(R_j=P_j(\sigma\alpha)/P_j(\alpha)\).
Thus \(R_j\in M^\times\) and \(v_i(R_j)=-\delta_{ij}\).

Choose an actual value \(b_\alpha\in(Mk)^\times\) of (F.49.7).
Since \(k/D\) is algebraic, a single finite extension
\(S_0/D\) inside \(k\) contains the coefficients needed to express
this element. The field
\[
                              F'=MS_0
\]
contains \(b_\alpha\), \(\sigma(b_\alpha)\), and every \(R_j\).
It is fixed once and for all, independently of the Kummer level \(n\).
It is \(\sigma\)-stable, because \(\sigma\) preserves \(M\) and
fixes \(S_0\subset k\).

If \(a\in(F')^\times\) and \(z^n=a\) with \(z\in Mk\), then
\(\sigma(z)/z\in F'\). Indeed, for
\(\tau\in\operatorname{Gal}(Mk/F')\), the automorphisms
\(\sigma,\tau\) commute on \(M\) and \(k\); \(\sigma\) fixes
\(\tau(z)/z\in\mu_n\). Thus every \(\tau\) fixes
\(\sigma(z)/z\). Consequently
\[
                      \sigma(a)/a\in((F')^\times)^n.
 \tag{F.49.9}
\]
Represent \(\lambda_j\bmod n\) by integers \(a_{j,n}\).
At level \(n\), evaluation of (F.49.7) says that the appropriate
quotient of \(b_\alpha\) by
\(b_0\prod_j(P_j(\alpha)/P_j(x_0))^{a_{j,n}}\)
is an \(n\)-th power in \(Mk\), and belongs to \(F'\).
Apply (F.49.9). The constants fixed by \(\sigma\) cancel, leaving
\[
 \frac{\sigma b_\alpha}{b_\alpha}
        \prod_j R_j^{-a_{j,n}}\in((F')^\times)^n.
 \tag{F.49.10}
\]
The integer-valued homomorphism \(v_i\circ N_{F'/M}\) gives, at
every level \(n\),
\[
 [F':M]\lambda_i
      =-(v_i\circ N_{F'/M})(\sigma b_\alpha/b_\alpha)
                  \in\mathbf Z
      \quad\hbox{inside }\widehat{\mathbf Z}.
 \tag{F.49.11}
\]
If a nonzero ordinary integer multiple of an element of
\(\widehat{\mathbf Z}\) is an ordinary integer, that element is
an ordinary integer: the corresponding rational quotient is
integral at every prime. Hence every \(\lambda_i\in\mathbf Z\).
This is the fixed-field norm argument of F.46, using normality,
roots of unity, and finite actual values. It uses neither a
pro-formation quotient nor divisibility of \(k^\times\).

\proofparagraph{From residues to actual units}
Let \(X/L\) be regular, connected and affine, with
\(x_0\in X(L)\), and let its pinned relative topos morphism to
\(\mathbb G_m\) have compatible Kummer class \(\xi\).
These are the source hypotheses of CHW Theorem~4.18(5).
In characteristic zero \(X\) is smooth; connectedness and a
rational point imply geometric connectedness, as in F.31.
Descend the geometry and generalized Albanese \(A\) to \(F\)
in (F.49.4), after the finite extension splitting the boundary
used in F.29.

The integral comparison of
\citet[Proposition~4.1, Corollary~4.2 and Proposition~4.3]{F_SpiessSzamuely2003}
identifies geometric \(H^1\) with the dual of the full Albanese
Tate module. Write \(H=G_L\triangleleft G_F\) and
\(W=(V_\ell A)_H\).
The projected point-Kummer images span \(W\).
For if a functional annihilated all of them, scale it to an
integral functional and choose a finite coefficient level at
which it is nonzero. It defines the finite arithmetic quotient
used in F.43--F.46's annihilating-functional argument;
(F.49.6) supplies a decomposition group surjecting onto that
quotient, contradicting the assumed vanishing.
The arithmetic quotient, rather than only its geometric kernel,
is detected here.

Choose finitely many spanning point images. Their coordinates
and actual multiplicative values descend to one finitely
generated field \(M/\mathbf Q\), enlarging \(F\) by finite data.
Normality of \(G_L\) makes \(W\) stable under \(G_F\).
Conjugating a point cocycle changes it by a coboundary that is
zero in \(H\)-coinvariants; its actual multiplicative value has
the cyclotomic transformation rule. The functional is therefore
\(G_M\)-equivariant on the chosen spanning images, hence on \(W\).
Apply \citet[Theorem~4.3(ii), p.~676]{F_Jannsen1995} over the
finitely generated field \(M\). It places this functional in the
integral unit lattice, for each coefficient prime, with the
boundary split and normalization at \(x_0\) as in F.29--F.30.
There is no appeal to a Tate-Hom theorem over the infinite
extension \(k\) or over \(L\).

For a curve, norm to a punctured rational line. The norm residue
is multiplied by the residue-field degree, as proved in F.5,
and (F.49.11) makes the original residue integral.
For a higher-dimensional source, take a smooth transverse curve
through a smooth point of each boundary divisor, after the finite
étale localization in F.27 and F.30. At a finite Kummer level the
local expression is \(u z^a\); transversality gives the same
coefficient on the curve. The punctured-line argument again
makes it integral. The injective divisor matrix of a unit basis
then yields an actual global unit \(h\).
Normalize \(h(x_0)=b_0\), the actual value at \(x_0\).
Its conjugates have equal divisors and ratio one at \(x_0\),
so it descends. The residual class is geometrically zero at
each finite Kummer level. Hochschild--Serre makes it constant,
and evaluation at \(x_0\) makes that constant zero. Thus
\[
                              \xi=\kappa(h).
 \tag{F.49.12}
\]
These are exactly the proved norm, transverse-curve and unit
descent arguments of F.27 and F.45--F.46, with their integrality
hypothesis supplied by (F.49.11). Equality (F.49.12) concerns the
compatible Kummer class; the remaining point comparison is needed
before the full-topos theorem can be applied.

\proofparagraph{Removing the scalar ambiguity in point values}
Let \(\varphi\) denote the actual geometric point function.
For a point over any finite extension, (F.49.5) and (F.49.12) give
\[
                         \varphi(x)=c_x h(x),\qquad c_x\in k^\times.
 \tag{F.49.13}
\]
Suppose first that the source is a curve and \(h\) is nonconstant.
If \(h(x)\in L\), the residue degree \([L(x):L]\) is at most the
degree of the finite map defined by \(h\). Thus sufficiently
high-degree points have \(h(x)\notin L\), and in particular
\(h(x)\notin k\). Their images in (F.49.13) are outside \(L\).
The point function is nonconstant, and the inverse image of
\(\mathbb G_m\setminus\{1\}\) is a nonempty open.

Pull back \(T-1\) there. The preceding argument gives a unit \(h_1\)
and \(d_x\in k^\times\) such that
\(\varphi(x)-1=d_xh_1(x)\).
Choose a point in this open of degree exceeding the total pole
degrees of \(h\) and \(h_1\) on the fixed smooth compactification.
The function
\[
                         c_xh-d_xh_1-1
\]
vanishes at that point; a nonzero function with that bounded
pole divisor cannot do so. It is identically zero.
Consequently \(h_1=Ah+B\), with \(A,B\in k^\times\).
At every point where \(h(x)\notin k\), comparison of the two
values in this identity gives \(c_x=-A/B=:c\).

For each \(b\in k^\times\), repeat on the nonempty inverse image
of \(\mathbb G_m\setminus\{b\}\). On high-degree points the
already determined coefficient is \(c\); the same pole-degree
argument gives a unit
\[
                              h_b=e_b(ch-b),
                       \qquad e_b\in k^\times.
\]
At a point with \(h(x)\in k\), choose \(b=ch(x)\).
If \(\varphi(x)\ne b\), this point belongs to the invertible
locus of the pullback of \(T-b\), but \(h_b(x)=0\), a contradiction.
Hence \(\varphi=ch\) at all points. The normalization at \(x_0\)
forces \(c=1\).

If \(h\) is constant, take a nontrivial finite extension \(L'/L\)
and choose \(\beta\in L'\setminus L\). Hilbertianity supplies
such an extension. Then \(\beta/h\notin k\).
All the preceding arguments remain valid over \(L'\).
By (F.49.13), every point value of \(\varphi\) is in \(hk\subset L\),
so no closed point maps to \(\beta\). Its inverse image is closed.
Any nonempty closed subset of a finite-type \(L'\)-scheme has a
closed point, so that inverse image is empty.
The function \(T-\beta\) therefore pulls back to a unit on the
whole source. Its actual values lie in \(L'\); (F.49.5) shows that
the corresponding unit also has all its point values in \(L'\).
If that unit were nonconstant, a point of degree exceeding its
pole degree would contradict this. It is constant.
Normalize it at \(x_0\). The resulting identities are
\[
          \frac{\varphi(x)-\beta}{h-\beta}\in k,\qquad
          \frac{\varphi(x)}{h}\in k.
 \tag{F.49.14}
\]
Since \(\beta/h\notin k\), these two scalar relations force
\(\varphi(x)=h\).

For higher-dimensional \(X\), pass to a finite extension and join
an arbitrary geometric point to \(x_0\) by a smooth geometrically
connected curve, as in F.29--F.31. The normalized units on that
curve have the same Kummer class. Their ratio is infinitely
divisible, hence constant by the divisor argument, and is one
at \(x_0\). The curve comparison therefore gives the value of
the same \(h\) at every geometric point of \(X\).

We have proved precisely condition (5) of
\citet[Theorem~4.18]{F_CHW2024}, for every finite extension as well.
Regular characteristic-zero schemes equal their absolute weak
normalizations. That theorem now gives the complete all-scheme
closed-point-preserving topos reconstruction over \(L\).
The conclusion uses equality of point functions followed by
CHW; it does not identify a full topos with its fundamental group.

\proofparagraph{Small geometric groups supply Hilbertianity}
Open subgroups and quotients of a small profinite group are small.
An open subgroup of index \(n\) in an open subgroup \(H\subset G\)
has index \(n[G:H]\) in \(G\). Open subgroups of a quotient pull
back injectively to subgroups of the same index. These observations
prove the assertions; no statement about arbitrary closed subgroups
is being used.

Suppose \(\operatorname{Gal}(k/K_0)\) is small, where
\(K_0=D\overline{\mathbf Q}\).
For \(L=Fk\) and \(B=F\overline{\mathbf Q}\), the field \(B/K_0\)
is finitely generated. Its relative algebraic closure of \(K_0\)
is finite \citep[Lemma~9.26.11, Tag~030D]{F_StacksFields}.
Thus \(k\cap B\) is finite over \(K_0\), and restriction identifies
\[
 \operatorname{Gal}(L/B)
     \simeq\operatorname{Gal}(k/(k\cap B)).
 \tag{F.49.15}
\]
This is an open subgroup of the small geometric group.
The field \(B/\overline{\mathbf Q}\) is finitely generated of
positive transcendence degree, hence Hilbertian.
Small Galois extensions of Hilbertian fields are Hilbertian by
\citet[Theorem~4.4, author p.~9]{F_BarySoroker2008}.
The more general Hilbert-property base-change result is
\citet[Theorem~1.2 and Corollary~4.3, pp.~1, 5]{F_BarySorokerFehmPetersen2023}.
We give the finite-cover argument here to check its arithmetic
constant quotient and its required point degrees.

Descend a connected finite Galois cover \(Y\to X\), with \(X/L\)
a smooth geometrically connected curve, to \(B\), enlarging \(F\)
by finitely many data as needed. Let \(M/B\) be the finite Galois
constant field of the descended cover. Its connectedness after
base change to \(L\) gives \(M\cap L=B\).
Choose a Riemann--Roch function of degree \(d\) on \(X_B\) and
put \(D_0=d|G|\).
There are finitely many subextensions of \(L/B\) of degree at
most \(D_0\). Let \(S/B\) be the compositum of their normal
closures inside \(L\). It is finite Galois, and
\(M\cap S=B\). Hence \(Y_B\times_B S\) is connected.
Hilbert irreducibility over \(B\), applied to the entire cover of
\(\mathbf P^1\) of degree \(D_0[S:B]\), gives one residue field
of that full degree. The \(Y_B\)-residue field \(E\) has degree
\(D_0\) over \(B\) and is linearly disjoint from \(S\).

The extension \(E\cap L\) has degree at most \(D_0\), so it lies
in \(S\). Therefore \(E\cap L=B\).
Since \(L/B\) is Galois, \(E\) and \(L\) are linearly disjoint.
The point degree \(d\) and the residual cover degree \(|G|\)
survive base change to \(L\). This gives (F.49.6) and unbounded
detecting degrees. Taking \(X\) to be an open of \(\mathbf A^1\)
and the Galois closure of an arbitrary irreducible separable
polynomial over \(L(T)\) proves Hilbertianity directly.
Its finitely many coefficients descend after increasing \(F\).
In particular \(k\) is Hilbertian, so the core theorem applies.

\proofparagraph{Positive-transcendence subfields of the small class}
Let \(k_0\subset L\) have positive transcendence degree over
\(\mathbf Q\). Take a finite transcendence basis \(T\) of
\(k_0/\mathbf Q\), put \(A=\mathbf Q(T)\), and enlarge \(F\)
in (F.49.4) to contain \(T\).
Let \(R\) be the relative algebraic closure of \(A\) in \(F\).
Then \(R/A\) is finite and \(F/R\) is regular.
Let \(N/R\) be the possibly infinite normal closure of \(Rk_0/R\).
Regularity identifies its \(R\)-conjugates with the corresponding
\(F\)-conjugates. Since \(L/F\) is normal, all of them lie in \(L\);
thus \(N\subset L\).
Regularity is preserved by extension of constants, giving
\[
 N\overline{\mathbf Q}\cap F\overline{\mathbf Q}
                         =R\overline{\mathbf Q}.
 \tag{F.49.16}
\]
Restriction makes
\(\operatorname{Gal}(N\overline{\mathbf Q}/R\overline{\mathbf Q})\)
a quotient of \(\operatorname{Gal}(L/B)\), hence small by
(F.49.15). The extension \(N\overline{\mathbf Q}/R\) is algebraic
Galois and contains \(\overline{\mathbf Q}\), while \(R\) is
finitely generated over \(\mathbf Q\) of positive transcendence
degree. The small-group case and the core theorem apply to it.

Finally \(N\overline{\mathbf Q}/k_0\) is algebraic:
\(R/A\) is finite, \(\overline{\mathbf Q}/\mathbf Q\) is algebraic,
and \(A\subset k_0\).
F.32 descends reconstruction to \(k_0\).
This proof uses only open subgroups and quotients of small groups.
It neither asserts closure under arbitrary closed subgroups nor
requires closure under finite subdirect products.
No equivalent intrinsic normal-closure classification is claimed
for this sufficient embedding class.

\proofparagraph{The alternating-product example}
Direct differentiation and evaluation give
\[
 P'_m(X)=(2m+1)(X^m-1)^2,\qquad
 P_m(\zeta)=\frac{2m^2}{m+1}\zeta\quad(\zeta^m=1).
 \tag{F.49.17}
\]
The \(m\) finite critical points have distinct critical values,
with exact ramification index three; infinity has index
\(n=2m+1\). The degree-\(n\) polynomial cover is geometrically
connected. Its finite branch inertia generators are three-cycles,
and these generate its geometric monodromy: quotient by their
normal closure would give a cover of \(\mathbf P^1\) unramified
outside infinity, which is trivial by Riemann--Hurwitz in
characteristic zero. All these cycles are even.

A transitive subgroup generated by three-cycles is the alternating
group on its letters. Here is the elementary support calculation.
The supports form a connected hypergraph, because its components
are group orbits. Begin with one support triple and its \(A_3\).
Adding a triple meeting the old set in two vertices extends the
alternating group to the one new vertex: conjugating by the old
alternating group supplies the required three-cycles.
If a triple meets the old set in one vertex and adds two, conjugate
it by an old three-cycle moving that vertex. The resulting pair
of triples share the two new vertices and generate \(A_4\) on
those two new and two old vertices. This reduces to successive
one-new-vertex steps. Induction proves the assertion.
Therefore the geometric splitting group is \(A_{2m+1}\).

For \(m\ge2\), these nonabelian simple groups are pairwise
nonisomorphic. Their finite collections of Galois extensions
over \(K_0=\overline{\mathbf Q}(t)\) are linearly disjoint:
a nontrivial common quotient of one simple group and a product
of the previous distinct simple groups would force an isomorphism
with a previous factor, by the elementary Goursat argument.
This proves (F.49.2).
Each \(\Omega_m/D\) and \(D\overline{\mathbf Q}/D\) is Galois,
so their compositum \(k_{\mathrm{alt}}/D\) is Galois.

The product in (F.49.2) is small. An open subgroup of index \(d\)
defines a continuous coset representation into \(S_d\).
A simple factor whose order is greater than \(d!\) has trivial
image, since any nontrivial homomorphism of a simple group
is injective. Only finitely many factors remain.
Continuity kills the entire product of the other factors:
their finite-support subgroup is dense and has already been killed.
The resulting finite product has finitely many subgroups.
Conversely every open subgroup contains all but finitely many
whole factors, so every open subgroup has alternating quotients
of unbounded degree.

Suppose this field had an alternative finitely generated
algebraic model \(F_0\subset k_{\mathrm{alt}}\) whose geometric
normal-closure group was virtually pro-\(\mathcal C\) for an
allowed proper formation \(\mathcal C\) of F.46.
Put \(F_1=F_0\mathbf Q(t)\).
It is finite over both \(F_0\) and \(\mathbf Q(t)\), since all
these fields have transcendence degree one and
\(k_{\mathrm{alt}}\) is algebraic over each.
Finite base change preserves virtual pro-\(\mathcal C\).
Normality over \(\mathbf Q(t)\) makes
\(\operatorname{Gal}(k_{\mathrm{alt}}/F_1\overline{\mathbf Q})\)
a quotient of that base-changed normal-closure group.
It would be virtually pro-\(\mathcal C\), but it is an open
subgroup of (F.49.2). Each of its open subgroups has arbitrarily
large alternating quotients, whereas F.44's elementary
proper-formation argument excludes every sufficiently large
\(A_n\) from \(\mathcal C\). This is a contradiction.
Thus the example is outside the full F.46 sufficient class,
including changes of model.

\proofparagraph{The universal-tripod example}
Let \(U=\mathbf P^1_{\overline{\mathbf Q}}\setminus\{0,1,\infty\}\).
The connected finite étale cover fields of \(U\) are stable
under composita and normal closure, so their union
\(k_{\mathrm{tri}}/\overline{\mathbf Q}(t)\) is Galois.
Every automorphism over \(\mathbf Q(t)\) preserves the algebraic
constants and the rational branch set \(\{0,1,\infty\}\).
It permutes these unramified extensions. Hence
\(k_{\mathrm{tri}}/\mathbf Q(t)\) is also algebraic Galois.

Riemann existence, \citet[Exposé~XII, Theorem~5.1, p.~333]{F_SGA1},
and geometric base change,
\citet[Exposé~XIII, 4.5.3, p.~421]{F_SGA1}, identify the group
with the profinite completion of the fundamental group of the
complex tripod. All primes are allowed in characteristic zero.
The complex tripod retracts onto a wedge of two circles, giving
(F.49.3). A two-generated profinite group is small: its coset
representations into \(S_d\) are determined by the images of two
generators, and there are finitely many point stabilizers for
each such finite action. The small-group theorem therefore
applies to \(k_{\mathrm{tri}}\).

Every open subgroup \(V\subset\widehat F_2\) is free profinite of rank
\[
                         1+[\widehat F_2:V]\ge2.
 \tag{F.49.18}
\]
For completeness, take the finite cover of the two-petal rose
corresponding to \(V\cap F_2\). The Euler-characteristic calculation
gives that free rank, and the finite-index subgroup has the full
induced profinite topology. Its completion is \(V\).
This is the same Schreier calculation as F.11.
For every odd \(n\ge5\), the group \(A_n\) is generated by
\((1\,2\,\cdots\,n)\) and \((1\,2\,3)\); conjugation gives
the adjacent three-cycles, which generate \(A_n\).
Thus every \(V\), and each of its open subgroups, surjects onto
alternating groups of unbounded degree. It cannot be virtually
pro-\(\mathcal C\) for any proper formation allowed in F.46.

The obstruction again survives every change of finitely generated
algebraic model. If \(F_0\subset k_{\mathrm{tri}}\) were such
a model, let \(N/F_0\overline{\mathbf Q}\) be its geometric
normal closure and put
\[
                         B=F_0\overline{\mathbf Q}\,\mathbf Q(t).
\]
This subfield of \(k_{\mathrm{tri}}\) is finite over both
\(F_0\overline{\mathbf Q}\) and \(\overline{\mathbf Q}(t)\):
both are finitely generated one-variable fields over
\(\overline{\mathbf Q}\) inside a common algebraic extension.
The group \(\operatorname{Gal}(N/B)\) is open in the proposed
virtually pro-\(\mathcal C\) group. Since
\(k_{\mathrm{tri}}/\overline{\mathbf Q}(t)\) is Galois,
restriction surjects onto
\(\operatorname{Gal}(k_{\mathrm{tri}}/B)\), an open subgroup
of \(\widehat F_2\). Quotients preserve virtual pro-\(\mathcal C\);
(F.49.18) and the unbounded alternating quotients contradict it.

Finally, the covers adjoining roots of \(t\) or \(1-t\) ramify
only at the removed points, so all these roots are in
\(k_{\mathrm{tri}}\). At the ordinary point \(t=2\), every
finite subextension is unramified. A coherent prolongation of
this discrete valuation to the union still has value group
\(\mathbf Z\) and gives \(v(t-2)=1\). Therefore \(t-2\) is not
a square in \(k_{\mathrm{tri}}\).
This explicitly shows why (F.49.5) is an inclusion, and why
radical closedness must not be added to the hypotheses.

\proofparagraph{Source and scope}
The exact finite supplement
\path{code/F_small_geometric_checks.py} and its JSON output
check eleven displayed polynomial identities, two alternating
generating-pair examples, and three support-induction examples.
These checks supplement the all-degree proofs above; they do not
test infinite field extensions or reconstruction.

The Hilbertianity of small Galois extensions is a known input,
attributed above with the numbering of the primary versions
actually opened. The proof retains the arithmetic constant
field and all point degrees required for reconstruction.
The generalized Albanese and integral Tate-Hom inputs are used
only at the previously proved and cited loci.
Both strict examples are proved directly, not by assuming that
reconstruction is preserved by filtered unions.
The focused primary-literature search and version distinctions
are recorded in \path{audit/F_small_geometric_literature.md}.
No priority or necessary-and-sufficient field classification is
claimed; F.9 remains unresolved.
\end{proof}
